\section{Final Specifications and Requirements}
The system was designed and evaluated under the following specifications.

\textbf{Hardware and compute.} All per-modality training and fusion training was executed on Google Colab Pro using an NVIDIA L4 GPU with 24 GB VRAM and 83 TFLOPS FP16 performance. Local development and analysis were performed on a Windows PC with Intel i7 CPU and 32 GB RAM, and in earlier phases on an Apple M1 MacBook Air. The absence of local GPU compute necessitated a hybrid workflow in which code is authored locally and executed remotely, with all data and checkpoints persisted to Google Drive at \texttt{MyDrive/Thesis\_EAV/}.

\textbf{Data specifications.} The EAV dataset provides pickle files of pre-extracted features for each of 42 subjects across three modalities. Each audio pickle contains Log-Mel Spectrogram segments of shape \texttt{(400, time, 128)} per subject, resampled to 16 kHz and segmented into five-second clips. Each vision pickle contains 25 video frames per clip, optionally face-cropped to $224\times224$ pixels. Each EEG pickle contains five-second signal windows at 100 Hz after downsampling from 500 Hz, band-pass filtered to 0.5--45 Hz, across 30 channels. The train/test split uses a stratified-by-class scheme (\texttt{h\_idx=56}) yielding 280 train and 120 test trials per subject per modality.

\textbf{Software and libraries.} The implementation uses PyTorch 2.10, the Hugging Face Transformers library 4.4x, torchaudio, torchvision, librosa, NumPy, Pandas, and scikit-learn. The EAV repository provides base data loaders and per-modality trainers, modified to fix the bugs documented in Chapter 7. The fusion module is implemented in a standalone \texttt{fusion/} directory at the project root, with no dependency on EAV repository internals.

\textbf{Training specifications.} Per-modality epoch budgets are: AST audio, 5 frozen + 5 fine-tuning epochs; ViT vision, 3 frozen + 2 fine-tuning epochs; EEGNet, 350 epochs from scratch. Fusion module training uses 80 epochs, the AdamW optimiser, learning rate $10^{-3}$, cosine annealing, batch size 32, and best-test-epoch checkpointing. The full 42-subject rollout required approximately 12 hours of wall-clock time on the Colab L4 GPU.

\textbf{Code and reproducibility.} All code is publicly archived in the GitHub repository \texttt{shafinmuffinn/thesis\_p2}. A single path-configuration module, \texttt{paths.py}, reads all filesystem locations from environment variables, enabling the same code to run unchanged in both local and Colab environments. Per-subject, per-modality logits are saved as compressed NumPy archives, enabling all downstream fusion and coherence analysis to be reproduced from cached results without re-running the expensive per-modality training.

\section{Societal and Clinical Impact}
The direct societal impact of this research is through the methodological contribution to affective computing: a validated framework for detecting cross-modal affective incoherence from passively collected multimodal data. Three downstream application domains stand to benefit.

\textbf{Mental health monitoring.} Automated, non-intrusive detection of emotion suppression or affective dysregulation could provide clinicians with objective biomarkers supplementing self-report instruments, which are known to be unreliable in populations with impaired introspective access such as depression, alexithymia, and dissociative disorders. The present work is a necessary precursor step---validating the methodology on healthy volunteers---before any clinical application can be responsibly pursued.

\textbf{Human-computer interaction.} Intelligent systems that can detect when a user's expressed affect diverges from their physiological state can adapt their behaviour accordingly: offering pauses in cognitively demanding interfaces when internal arousal is high, or re-evaluating the affective valence of a recommendation when the user's EEG indicates a mismatch with their expressed response.

\textbf{Affective computing research.} By providing the first characterisation of cross-modal affective coherence patterns on the EAV benchmark, this thesis creates a baseline against which future, more sophisticated coherence-detection methods can be compared.

\section{Ethical Considerations}
The EAV dataset is an ethically approved corpus collected by Lee, Shomanov, Kabidenova, Yazici, and colleagues at Nazarbayev University. Participants provided informed consent to recording and to the release of de-identified data for research purposes. The dataset is not openly distributed; access is granted on individual request to the dataset authors. For this thesis, access was granted by Adai Shomanov via a token-authorised Zenodo link, with explicit permission for the dataset's use in this research. This thesis uses the pre-processed EAV pickle files in exactly the conditions intended by the data providers; no additional data collection or human-participants research was conducted.

The suppression matrix methodology, applied to clinical populations in future work, raises ethical considerations around privacy, stigmatisation, and clinical deployment without appropriate validation. These are explicitly out of scope for the present thesis, which is confined to analysis of laboratory data. Any future work involving clinical populations must obtain independent ethics approval and must not make clinical inferences based solely on the present thesis's findings, which have not been validated against clinical ground truth.
